\documentclass[11pt]{article}

\usepackage{../ppackage}
\usepackage{bbm}
\usepackage{import}






\usepackage{tikz}
\usetikzlibrary{arrows.meta,patterns}
\usetikzlibrary{ipe} % ipe compatibility library

\usepackage{../../../Notes/tikzit}
\usepackage{../../../Notes/utf8math}

\input{../../../Notes/TIKZ/digraph.tikzstyles}



\usepackage{url}


\newcommand{\solution}{
\bigskip\noindent
	\textbf{Solution: \\}}
	


\DeclareMathOperator{\size}{size}
\DeclareMathOperator{\conv}{conv}
\newcommand{\SV}{\mathrm{SV}}
\newcommand{\bigO}{O}
\newcommand{\cut}{\mathrm{cut}}
\newcommand{\LLL}{\mathrm{LLL}}
\newcommand{\setR}{\mathbb{R}}
\newcommand{\setZ}{\mathbb{Z}}
\newcommand{\setQ}{\mathbb{Q}}
\newcommand{\setC}{\mathbb{C}}
\newcommand{\setN}{\mathbb{N}}
\newcommand{\wt}[1]{\widetilde{#1}}
\newcommand{\opt}{{\sc 0/1-opt}\xspace}
\newcommand{\aug}{{\sc 0/1-aug}\xspace}
\newcommand{\psep}{{\sc 0/1-psep}\xspace}
\newcommand{\sep}{{\sc 0/1-sep}\xspace}
\newcommand{\fopt}{{\sc 0/1-testopt\xspace} }

\newcommand{\hpp}{\mathrm{HPP}}
\newcommand{\nodes}{\mathcal{V}}
\newcommand{\vol}{\mathrm{vol}}
\newcommand{\diag}{\mathrm{diag}}
\newcommand{\arcs}{\mathcal{A}}
\newcommand{\edges}{\mathcal{E}}
\newcommand{\paths}{\mathscr{P}}
\newcommand{\cycles}{\mathcal{C}}




\newcommand{\K}{{\mathcal K}}
\newcommand{\A}{{A}}
\newcommand{\B}{{B}}
\newcommand{\T}{\mathscr{T}}
\newcommand{\eE}{\mathscr{E}}
\newcommand{\eS}{\mathscr{S}}
\newcommand{\eP}{\mathscr{P}}
\newcommand{\eM}{\mathscr{M}}



\newcommand{\transp}{^{\mathrm{T}}}

\newcommand{\smallmat}[1]{\left( \begin{smallmatrix} #1 \end{smallmatrix}\right)}

\newcommand{\mat}[1]{ \begin{pmatrix} #1 \end{pmatrix}}
\newcommand{\smat}[1]{ \big(\begin{smallmatrix} #1 \end{smallmatrix}\big)}

\newcommand{\pc}{\mathscr{P}}
\newcommand{\ob}{\mathscr{O}}
\newcommand{\odds}{\mathscr{W}}
\newcommand{\up}{\mathscr{U}}
\newcommand{\ef}{\mathscr{F}}
\newcommand{\eh}{\mathscr{H}}
\newcommand{\ev}{\mathscr{V}}
\newcommand{\ec}{\mathscr{C}}
\newcommand{\eu}{\mathscr{U}}

\newcommand{\lex}{\mathrm{lex}}

\renewcommand{\leq}{\leqslant}
\renewcommand{\geq}{\geqslant}









\newcommand{\linhull}{\mathrm{lin.hull}}
\newcommand{\affhull}{\mathrm{affine.hull}}
\newcommand{\charcone}{\mathrm{char.cone}}
\newcommand{\cone}{\mathrm{cone}}
\newcommand{\rank}{\mathrm{rank}}
\newcommand{\wb}[1]{\overline{#1}}



\usepackage{enumerate}

      
\institute{\'Ecole Polytechnique F\'ed\'erale de Lausanne}
\lecture{Discrete Optimization}
\faculty{Prof. Friedrich Eisenbrand}
\term{Spring 2026}
\publishdate{March 24, 2026}
\duedate{ }
\problemset{Problem Set~5 (Solutions)}

\begin{document}
\makeheader

\begin{enumerate}[1)]

\item Let $P= \{x ∈\setR^n : Ax ≤b\}$ be a bounded, non-empty polyhedron. Formulate a linear program that
computes the largest Euclidean ball inside P.

\begin{solution}
A ball of radius $r$ and center $x$ is contained in $P$ if and only if $x ∈P$ and $x$ has distance at least $r$
from any hyperplane defining $P$. Hence we obtain the following linear program:
\begin{align*}
\max \quad &r \\
\text{subject to }& \frac{b_i−a_i x}{\Vert a_i \Vert} ≥r \quad \forall i = 1, \hdots, m \\
& Ax ≤b
\end{align*}
where $a_1, \hdots , a_m$ are the rows of $A$ and $b = (b_1 . . . b_m)^T$.
\end{solution}

\item Suppose you are given an oracle algorithm, which for a given polyhedron
$$P= \{\bar{x}∈\setR^n: \bar{A}\bar{x} \leq \bar{b}\}$$
gives you a feasible solution or asserts that there is none. Show that using a single call of this oracle
one can obtain an optimum solution for the LP
$$\max\{c^T x : x∈\setR^n; Ax≤b\}$$
assuming that the LP is feasible and bounded.

\begin{solution}
The LP is feasible and bounded, thus an optimum solution must exist. Strong duality tells us that
the dual $\min\{b^T y: A^T y= c, y≥0\}$ is feasible and bounded. For optimal solutions $x^∗$ of the primal
and $y^∗$ of the dual we have $$b^T y^∗= c^T x^∗.$$

Thus every point $(x^∗,y^∗)$ of the polyhedron
\begin{align*}
&c^T x = b^T y \\
& Ax ≤ b \\
&A^T y= c \\
& y ≥ 0
\end{align*}
is optimal. Hence with one oracle call for the polyhedron above we get an optimal solution of the
LP.
\end{solution}


% \item Determine the dual program for the following linear program:
% \begin{align*}
% \min \quad& 3x_1 + 2x_2−3x_3 + 4x_4 \\
% &2x_1−2x_2 + 3x_3 + 4x_4 ≤ 3 \\
% & x2 + 3x3 + 4x4 ≥ −5 \\
% &2x_1−3x_2−7x_3−4x_4 = 2 \\
% &x_1 ≥0 \\
% & x_4 ≤0
% \end{align*}

% \begin{solution}
% \begin{align*}
% \max \quad &3y_1−5y_2 + 2y_3 \\
% & 2y_1 + 2y_3 ≤ 3 \\
% &−2y_1 + y_2−3y_3 = 2 \\
% &3y_1 + 3y_2−7y_3 =−3 \\
% & y_1 + y_2−y_3 ≥ 1\\
% &y_1 ≤0 \\
% & y_2 ≥0
% \end{align*}
% \end{solution}


\item Consider a linear program $\max\{c^Tx: x\in \R^n, Ax\leq b\}$ where $A\in \R^{m\times n}$ is not full column rank. 
\begin{enumerate}[i)]
\item Convert this LP to an equivalent LP with a full column rank constraint matrix. 

\emph{Hint: Write $x$ as difference of two non-negative vectors.}

\item Write down the dual of the converted LP in i).
\item Construct a feasible solution to dual of the original LP, based on a given feasible solution to the dual of the converted LP in ii).
\end{enumerate}
One can conclude from the above that the strong duality theorem holds when the constraint matrix of LP does not have full column rank (assuming it holds for full column rank constraint matrices).

\begin{solution}
\begin{enumerate}[i)]
\item Write $x = x^+ - x^-$ where $x^+, x^- \geq 0$. Then the original LP is equivalent to
\begin{align*}
\max \quad & c^T x^+ - c^T x^- \\
\text{subject to }& A x^+ - A x^- ≤ b \\
& x^+ ≥ 0 \\
& x^- ≥ 0
\end{align*}
where the constraint matrix is of full column rank.
\item The dual of the converted LP in i) is
\begin{align*}
\min \quad & (b\;\,0\;\,0)^T \begin{pmatrix}
y_1 \\ y_2 \\ y_3	
\end{pmatrix} \\
\text{subject to }& y_1^T A - y_2^T = c^T \\
& -y_1^T A - y_3^T = -c^T \\
& y_1 \in \R^m, y_2, y_3\in \R^n, y_1, y_2, y_3 ≥ 0
\end{align*}
\item Let $(y_1, y_2, y_3)^T$ be a feasible solution to the dual of the converted LP. By adding up the two constraints we get $y_2+y_3=0$. Since $y_2, y_3 ≥ 0$, we have $y_2 = y_3 = 0$. Thus $y_1^T A = c^T$ and $y_1$ is a feasible solution to the dual of the original LP.
\end{enumerate}
\end{solution}

\item Consider the following linear program:
\begin{align*}
\max \quad & x_1 + x_2 \\
\text{subject to }&2x_1 + x_2 ≤6 \\
&x_1 + 2x_2 ≤8 \\
& 3x_1 + 4x_2 ≤ 22\\
& x_1 + 5x_2 ≤ 23
\end{align*}
Show that $(4/3, 10/3)$ is an optimal solution by using weak duality.


\begin{solution}
The assignment $(4/3, 10/3)$ has the objective function value of $14/3$. In order to prove that it is
optimal (via weak duality), we are going to form the dual LP, and find a feasible solution to the
dual that achieves the same objective value. The dual is:
\begin{align}
\min  \quad & 6y_1 + 8y_2 + 22y_3 + 23y_4 \\
\text{subject to }& 2y_1 + y_2 + 3y_3 + y_4 = 1 \\
&y_1 + 2y_2 + 4y_3 + 5y_4 = 1\\
 & y1, y2, y3, y4 ≥0
\end{align}
Thus, we are looking for a feasible dual solution such that $6y_1 + 8y_2 + 22y_3 + 23y_4 = 14/3$. By using
Gaussian elimination on this constraint combined with (1) and (2) we get:
$$−4y_2−2y_3−7y_4 =−4/3$$
$$−3y_2−5y_3−9y_4 =−1$$
and 
$$14/3y_3 + 5y_4 = 0.$$
Since $y_1, y_2, y_3, y_4 ≥0$, we have that $y_3 = y_4 = 0$ and then $y_1 = y_2 = 1/3$. This is the desired
feasible dual solution coinciding with the primal solution $(4/3, 10/3)$, proving the optimality of the
latter.

\end{solution}


\item In this exercise, we will use LP duality to prove Farkas' lemma: for any $A\in \R^{m\times n}$ and $b\in \R^m$, $Ax \leq b$ is feasible if and only if for any $\lambda \in \R^m_{\geq 0}$, $\lambda^T A = 0$ implies that $\lambda^T b \geq 0$.
For the if direction, do the following steps. 

Let $\{Ax\leq b\}$ be feasible and $\{Ax\leq b, \alpha^T x\leq \beta\}$ be infeasible.
\begin{enumerate}[i)]
	\item Show that the LP $\max\{(-\alpha)^Tx: Ax\leq b\}$ is bounded. Write down its dual.
	\item Using strong duality, show that there exists $\lambda \in \R^m_{\geq 0}$ such that $(\lambda \;\, 1)^T \begin{pmatrix} A \\ \alpha^T \end{pmatrix} = 0$ and $(\lambda \;\, 1)^T \begin{pmatrix} b \\ \beta \end{pmatrix} < 0$.
	\item Using ii), conclude the proof of Farkas' lemma by induction on $m$.
\end{enumerate}

\begin{solution}
	\begin{enumerate}[i)]
		\item Since $\{Ax\leq b\}$ is feasible and $\{Ax\leq b, \alpha^T x\leq \beta\}$ is infeasible, for any $x$ such that $Ax\leq b$ we have $\alpha^T x > \beta$, i.e., $(-\alpha)^T x < -\beta$. 
		Its dual is $\min\{b^T \lambda: \lambda^T A = -\alpha^T, \lambda≥0\}$.
		\item By strong duality, there exists $\lambda \in \R^m_{\geq 0}$ such that $\lambda^T A = -\alpha^T$ and $\lambda^T b < -\beta$. Augment $\lambda$ with $1$ to get what we need.
		\item Finally we prove the if direction of Farkas' lemma by induction on the number of constraints $m$.
		If $m=0$, the statement is trivial. 
		Assume that the statement holds for $\leq m$. Consider ${A^\prime x \leq b^\prime}$ infeasible where $A^\prime \in \R^{(m+1)\times n}$ and $b^\prime \in \R^{m+1}$. Let $Ax\leq b$ be the first $m$ constraints and $\alpha^T x \leq \beta$ be the last constraint of $A^\prime x \leq b^\prime$.  
		If $Ax \leq b$ is infeasible, by induction hypothesis there exists $\lambda \in \R^m_{\geq 0}$ such that $\lambda^T A = 0$ and $\lambda^T b < 0$. Augment $\lambda$ with $0$ to get what we need. 
		If $Ax \leq b$ is feasible, by ii) there exists $\lambda \in \R^m_{\geq 0}$ such that $(\lambda \;\, 1)^T \begin{pmatrix} A \\ \alpha^T \end{pmatrix} = 0$ and $(\lambda \;\, 1)^T \begin{pmatrix} b \\ \beta \end{pmatrix} < 0$. This concludes the proof.
	\end{enumerate}
\end{solution}


\end{enumerate}



  

\end{document}

%%% Local Variables:
%%% mode: latex
%%% TeX-master: t
%%% End:
